% ============================================================
\chapter{Quadrature de Gauss}
\label{ch:gauss}
% ============================================================

\section{Exemple motivant}
Les formules de Newton--Cotes (n\oe uds fix\'es) atteignent un degr\'e d'exactitude $n$ ou $n+1$. En choisissant \textbf{les n\oe uds librement}, on peut atteindre le degr\'e $2n+1$ avec $n+1$ points --- c'est la quadrature de Gauss.

\section{Principe}

\begin{definition}[Quadrature de Gauss]\index{quadrature de Gauss}
Trouver $n+1$ n\oe uds $x_0,\ldots,x_n$ et poids $w_0,\ldots,w_n$ tels que :
\[
\int_{-1}^1 f(x)\,dx \approx \sum_{i=0}^n w_i f(x_i)
\]
soit exacte pour tout $f\in\mathcal{P}_{2n+1}$.
\end{definition}

\begin{theoreme}[Th\'eor\`eme fondamental]
Les n\oe uds de la quadrature de Gauss--Legendre \`a $n+1$ points sont les racines du polyn\^ome de Legendre $P_{n+1}$. Le degr\'e d'exactitude est $2n+1$.
\end{theoreme}

\begin{proof}[Esquisse]
Soit $f\in\mathcal{P}_{2n+1}$. Division euclidienne : $f=q\cdot P_{n+1}+r$ avec $\deg q\leq n$, $\deg r\leq n$. Alors $\int_{-1}^1 f=\int_{-1}^1 qP_{n+1}+\int_{-1}^1 r$. Par orthogonalit\'e de $P_{n+1}$ \`a $\mathcal{P}_n$ : $\int qP_{n+1}=0$. Et $\sum w_i f(x_i)=\sum w_i r(x_i)$ car $P_{n+1}(x_i)=0$. Les poids sont choisis pour que la formule soit exacte sur $\mathcal{P}_n$, donc $\sum w_i r(x_i)=\int r$.
\end{proof}

\section{N\oe uds et poids}

\begin{center}
\begin{tabular}{cll}
\toprule
$n+1$ & N\oe uds $x_i$ & Poids $w_i$ \\
\midrule
1 & $0$ & $2$ \\
2 & $\pm 1/\sqrt{3}$ & $1,\;1$ \\
3 & $0,\;\pm\sqrt{3/5}$ & $8/9,\;5/9,\;5/9$ \\
\bottomrule
\end{tabular}
\end{center}

\section{Gauss--Legendre, Gauss--Laguerre, Gauss--Hermite}

\begin{definition}
Pour $\int_a^b f(x)w(x)\,dx$ avec poids $w$ :
\begin{itemize}
\item \textbf{Gauss--Legendre} : $w(x)=1$ sur $[-1,1]$.
\item \textbf{Gauss--Laguerre} : $w(x)=e^{-x}$ sur $[0,\infty)$.
\item \textbf{Gauss--Hermite} : $w(x)=e^{-x^2}$ sur $(-\infty,\infty)$.
\end{itemize}
\end{definition}

\section{Erreur}

\begin{theoreme}
L'erreur de Gauss--Legendre \`a $n+1$ points est :
\[
E_n(f) = \frac{2^{2n+3}[(n+1)!]^4}{(2n+3)[(2n+2)!]^3}f^{(2n+2)}(\xi).
\]
\end{theoreme}

\section{Changement d'intervalle}

Pour $\int_a^b f(x)\,dx$, le changement $x=\frac{b-a}{2}t+\frac{a+b}{2}$ ram\`ene \`a $[-1,1]$ :
\[
\int_a^b f(x)\,dx = \frac{b-a}{2}\int_{-1}^1 f\!\left(\frac{b-a}{2}t+\frac{a+b}{2}\right)dt.
\]

\section{Impl\'ementation}

\begin{codeblock}[title=\textbf{Python}]
\begin{minted}{python}
import numpy as np
from numpy.polynomial.legendre import leggauss

def gauss_legendre(f, a, b, n):
    """Quadrature de Gauss-Legendre a n points sur [a,b]."""
    nodes, weights = leggauss(n)
    # Changement d'intervalle [-1,1] -> [a,b]
    x = 0.5*(b-a)*nodes + 0.5*(a+b)
    w = 0.5*(b-a)*weights
    return np.sum(w * f(x))

f = lambda x: np.exp(-x**2)
exact = 0.7468241328124271

for n in [2, 3, 5, 10]:
    approx = gauss_legendre(f, 0, 1, n)
    print(f"n={n:2d}: I={approx:.15f}, erreur={abs(approx-exact):.2e}")
\end{minted}
\end{codeblock}

\begin{codeblock}[title=\textbf{Julia}]
\begin{minted}{julia}
using FastGaussQuadrature

function gauss_legendre_ab(f, a, b, n)
    nodes, weights = gausslegendre(n)
    x = 0.5(b-a) .* nodes .+ 0.5(a+b)
    w = 0.5(b-a) .* weights
    return sum(w .* f.(x))
end

f(x) = exp(-x^2)
for n in [2, 3, 5, 10]
    I = gauss_legendre_ab(f, 0, 1, n)
    println("n=$n: I=$I")
end
\end{minted}
\end{codeblock}

\section{Exercices}

\begin{exercice}[$\star$]
Calculer $\int_0^1 x^3\,dx$ par Gauss--Legendre \`a 2 points. V\'erifier l'exactitude.
\end{exercice}

\begin{exercice}[$\star\star$]
Montrer que les poids de Gauss sont toujours positifs.
\end{exercice}

\begin{exercice}[$\star\star\star$ -- Projet]
Comparer la convergence de Simpson composite ($n$ points) et Gauss--Legendre ($n$ points) pour $\int_0^1 e^{-x^2}dx$. Tracer l'erreur en \'echelle log-log.
\end{exercice}

\begin{formules}
\begin{align*}
\int_{-1}^1 f(x)\,dx &\approx \sum_{i=0}^n w_i f(x_i), \quad \text{exacte pour }\mathcal{P}_{2n+1} \\
\int_a^b f(x)\,dx &= \frac{b-a}{2}\int_{-1}^1 f\!\left(\frac{b-a}{2}t+\frac{a+b}{2}\right)dt
\end{align*}
\end{formules}
